\section{Summary}

\noindent The 2D XY model has Hamiltonian $$H = -J \sum_{\langle i,j \rangle} \cos(\theta_i - \theta_j)$$,
with $J=1$ and spins $\theta_i \in [0, 2\pi)$ shows the Kosterlitz-Thouless (KT) transition at $T_c \approx 0.893$. Below $T_c$, bound vortex-antivortex pairs maintain quasi-long-range order; above $T_c$, they unbind, causing disorder.

\noindent Monte Carlo simulations compare Metropolis single-spin flips and Wolff cluster flips on lattices $L=8$ to $256$, with $T \in [0.1, 1.5]$ in steps of $0.02$.
Susceptibility peaks sharpen and shift to $T_c$ as $L$ grows. Metropolis suffers strong critical slowing down ($\tau_{\texttt{exp}} \sim 10^3-10^5$ near $T_c$), while Wolff reduces it by 1-2 orders ($\tau \sim 10-10^3$). Meaning in order to get accurate results for the Metropolis algorithm, you would need to take significantly more measurements than for the Wolff algorithm, which is more efficient in this regard.

\noindent Furthermore Finite-size scaling was implemented. This leads to susceptibility peaks giving $T_c = 0.884 \pm 0.026$, $\nu = 0.562 \pm 0.057$, matching known values $T_c = 0.893$, $\nu=0.5$.
A data collapse of $\chi(L,T)$ using $\eta \approx 0.01$, $b \approx 1.68$ confirms scaling, consistent with KT theory predicting $\eta(T_c)=1/4$.

\noindent Wolff excels near transitions in symmetric models, unlike local Metropolis updates. Metropolis works for trends, but clusters are needed for precise $T_c$ and exponents. Overall the superior performance of Wolff algorithm in reducing autocorrelation times and improving accuracy near criticality is evident, making it the preferred choice for studying the XY model's KT transition.